% -*- coding: utf-8 -*-
% !TEX program = xelatex
\documentclass[plain,noanswer,medmath]{../../template/jnuexam}
\usepackage{../../template/kysx}

\begin{document}


\examtitle{2012}{2}


\loadproblems{1}{2012P1}%载入试卷一


\makepart{选择题}{1～8小题，每小题4分，共32分}


\useproblem{1}{1}{1}%【同试卷一第一(1)题】


\useproblem{1}{1}{2}%【同试卷一第一(2)题】


\begin{problem}
设$a_n > 0$($n = 1,2, \cdots$)，$S_n = a_1 + a_2 + \cdots + a_n$，
则数列$\{s_n\}$有界是数列$\{a_n\}$收敛的\pickout{}
\begin{abcd}
\item 充分必要条件．
\item 充分非必要条件．
\item 必要非充分条件．
\item 即非充分也非必要条件．
\end{abcd}
\end{problem}


\useproblem{1}{1}{4}%【同试卷一第一(4)题】


\begin{problem}
设函数$f\left(x,y \right)$可微，且对任意$x,y$都有
$\frac{\pd f(x,y)}{\pdx} > 0$，$\frac{\pd f(x,y)}{\pdy} < 0$．
则使不等式$f\left(x_1,y_1 \right) < f\left(x_2,y_2 \right)$成立的一个充分条件是
\begin{abcd}
\item $x_1 > x_2,y_1 < y_2$．
\item $x_1 > x_2,y_1 > y_1$．
\item $x_1 < x_2,y_1 < y_2$．
\item $x_1 < x_2,y_1 > y_2$．
\end{abcd}
\end{problem}


\begin{problem}
设区域D由曲线$y = \sin x,x = \pm \frac{\pi}{2},y = 1$围成，则$\iint_D (x^5 y - 1)\dx\dy =$\pickout{}
\begin{abcd}
\item $\pi$．
\item $2$．
\item $- 2$．
\item $- \pi$．
\end{abcd}
\end{problem}


\useproblem{1}{1}{5}%【同试卷一第一(5)题】


\useproblem{1}{1}{6}%【同试卷一第一(6)题】


\makepart{填空题}{9～14小题，每小题4分，共24分}


\begin{problem}
设$y = y(x)$是由方程$x^2 - y + 1 = \e^y$所确定的隐函数，
则$\left. \frac{\d^2y}{\dx^2} \right|_{x = 0} =$ \fillin{}．
\end{problem}


\begin{problem}
计算$\lim_{x \to \infty} n\left(\frac{1}{1 + n^2} + \frac{1}{2^2 + n^2}
+ \cdots + \frac{1}{n^2 + n^2} \right) =$ \fillin{}．
\end{problem}


\begin{problem}
设$z = f\left(\ln x + \frac{1}{y} \right)$，其中函数$f(u)$可微，则$x\frac{\pdz}{\pdx} + y^2\frac{\pdz}{\pdy} =$ \fillin{}．
\end{problem}


\begin{problem}
微分方程$y\dx + (x - 3y^2)\dy = 0$满足初始条件$y|_{x = 1}$=1的解为$y =$ \fillin{}．
\end{problem}


\begin{problem}
曲线$y = x^2 + x$（$x < 0$）上曲率为$\frac{\sqrt{2}}{2}$的点的坐标是 \fillin{}．
\end{problem}


\begin{problem}
设$A$为3阶矩阵，$|A| = 3$，$A^*$为$A$的伴随矩阵，
若交换$A$的第一行与第二行得到矩阵$B$，则$\left| BA^* \right| =$ \fillin{}．
\end{problem}


\makepart{解答题}{15～23小题，共94分}


\begin{problem}[points=10]
已知函数$f(x) = \frac{1 + x}{\sin x} - \frac{1}{x,}$，记$a = \lim_{x \to 0} f(x)$．\par
\begin{enumerate*}
\item 求$a$的值；
\item 若当$x \to 0$时，$f(x) - a$是$x^k$的同阶无穷小，求$k$．
\end{enumerate*}
\end{problem}


\useproblem[points=10]{1}{3}{16}%【同试卷一第三(16)题】


\begin{problem}[points=10]
过$(0,1)$点作曲线$L:\;y = \ln x$的切线，切点为$A$，又$L$与$x$轴交于$B$点，区域$D$由$L$与直线$AB$围成，
求区域$D$的面积及$D$绕$x$轴旋转一周所得旋转体的体积．
\end{problem}


\begin{problem}[points=10]计算二重积分$\iint_D xy\d\sigma$，其中区域D为曲线
$$r = 1 + \cos \theta \;\left(0 \le \theta \le \pi \right)$$
与极轴围成．
\end{problem}


\begin{problem}[points=11]
设函数$f(x)$满足方程$f''(x) + f'(x) - 2f(x) = 0$及$f'(x) + f(x) = 2\e^x$．
\begin{enumerate}
\item 求表达式$f(x)$；
\item 求曲线的拐点$y = f(x^2)\int_0^x f(- t^2)\dt$．
\end{enumerate}
\end{problem}


\useproblem[points=10]{1}{3}{15}%【同试卷一第三(15)题】


\begin{problem}[points=11]
\begin{enumerate}
\item
证明方程$x^n + x^{n - 1} + \cdots + x = 1$（$n$为大于$1$的整数）%
在区间$\left(\frac{1}{2},1 \right)$内有且仅有一个实根；
\item
记(I)中的实根为$x_n$，证明$\lim_{n \to \infty} x_n$存在，并求此极限．
\end{enumerate}
\end{problem}


\useproblem[points=11]{1}{3}{20}%【同试卷一第三(20)题】


\useproblem[points=11]{1}{3}{21}%【同试卷一第三(21)题】


\end{document}
